\section{Glossary}
\label{sec:glossary}

Each entry reads informally first, then formally. The formal definitions are
normative.

\begin{notationx}[Metavariables]\label{not:vars}
$P, Q$ range over processes; $p, q, r$ over ranks; $g$ over global element
indices; $a$ over addresses; $t$ over origin times, $\delta$ over deltas,
$\eff$ over effective times. Processes and ranks are never identified.
\end{notationx}

% ===========================================================================
\subsection{Identity and structure}
\label{sec:gloss-identity}

\term{process}{Process}
A process is one participant in a distributed simulation: a running instance
of the program, in practice one MPI process. Processes are the entities that
act --- they own memory, advance along the timeline and issue operations.
Their number is fixed for a run: DSTM performs no load balancing, spawning or
migration.

\term{rank}{Rank}
A rank is not a participant but a name for one: an integer identifying a
process uniquely. Ranks, being integers, can be ordered, indexed and arranged
in a ring; a process as such cannot. Every structure in this document is
indexed by rank.

\begin{definition}[Processes and ranks]\label{def:process}
Let $\worldsize \in \Nat$, $\worldsize \ge 1$, be the number of processes,
fixed for a \gls{session}. The \emph{rank set} is
\[
  \Ranks \defeq \{0, 1, \dots, \worldsize - 1\} \subset \Nat ,
\]
totally ordered, with \emph{ring successor} $\succr{p} \defeq (p+1) \bmod
\worldsize$. The \emph{processes} are the family
$(\proc{i})_{0 \le i < \worldsize}$ of pairwise distinct entities, with
$\Procs \defeq \{\proc{i} : i \in \Ranks\}$. The map
$\rankof{\proc{i}} \defeq i$ is a bijection $\Procs \to \Ranks$.
\end{definition}

\begin{remark}[The bijection is an assumption of this model]
\label{rem:rank-process}
$\rankof{\cdot}$ is bijective only because DSTM fixes the assignment for a
whole \gls{session}: no process is created, destroyed, migrated or
renumbered. A model admitting spawning, failure, migration, or several ranks
per process would need $\rankof{\cdot}$ partial, time-dependent or
non-injective, and every rank-indexed structure below would have to be
re-attached to processes. That generality --- logical processes, virtual
agents, load balancing --- is the extension \cite{dstm2026} defers, and
\cref{def:process} is where it would be made.
\end{remark}

\term{session}{Session}
A session is one interval of DSTM activity, opened by \code{dstm\_start} and
closed by \code{dstm\_fini}. Within it the \gls{comm-graph} and every region
size are fixed; changing either means closing the session and opening
another. A process runs sessions in sequence, never two at once.

\begin{definition}[Session]\label{def:session}
A \emph{session} is a maximal interval of execution during which $\cgraph$ and
every $S_p$ are constant. All processes enter and leave it collectively.
\end{definition}

\term{region}{Region}
\alsoterm{address-space}
A region is the flat, byte-addressable memory owned by one process, sized at
session start. Each process owns exactly one, named by that process's rank:
region identifiers and rank identifiers are the same numbers.

\begin{definition}[Region and address space]\label{def:region}
The \emph{local address space} of rank $p$ is
$\Addrof{p} \defeq \{0, \dots, S_p - 1\}$ for a region size $S_p > 0$; the
\emph{region} of rank $p$ is the corresponding partition of the global address
space $\Addr$. Both are fixed by \cref{ax:pgas}.
\end{definition}

\term{allocation}{Allocation}
Data structures carve reserved extents out of a region. An allocation is one
such reservation: a contiguous run of addresses set aside for one structure
and, while it lives, off limits to every other.

Two points are easy to get wrong. An allocation is a property of
\emph{addresses}, not \glspl{cell}: it does not vary with time, and an
allocated address is allocated in every \gls{version} alike. And what is
specified is the \emph{contract} an allocator meets, never the algorithm;
which addresses come back is deliberately open.

\begin{definition}[Allocation]\label{def:allocation}
Fix a set $\Handles$ of \emph{handles}. The allocation state of rank $p$ is a
partial function
\[
  \allocof{p} \;:\; \Addrof{p} \rightharpoonup \Handles ,
\]
mapping each reserved address to the handle reserving it. Write
$\live{p} \defeq \operatorname{dom}(\allocof{p})$ for the \emph{allocated}
addresses and $\extent{p}{h} \defeq \allocof{p}^{-1}(h)$ for the \emph{extent}
of a live handle. The state is well formed when every extent is a contiguous
interval of $\Addrof{p}$.
\end{definition}

Being a function, not a predicate, is what does the work: no address carries
two handles, so \emph{distinct allocations are disjoint automatically}. That
is what derived structures rely on, and why a handle-valued map beats a
boolean.

\begin{definition}[Allocation operations]\label{def:alloc-ops}
$\mathsf{alloc}_p(n)$ yields a fresh handle $h$ and an interval
$I \subseteq \Addrof{p}$ of length $n$ disjoint from $\live{p}$, extending
$\allocof{p}$ by $I \mapsto h$; it fails if no such interval exists.
$\mathsf{free}_p(h)$ removes $\extent{p}{h}$ from the domain. Between the two,
$\extent{p}{h}$ is constant: allocations never move or resize.
\end{definition}

\begin{remark}[The policy is not specified]\label{rem:alloc-policy}
\Cref{def:alloc-ops} says which intervals are admissible, not which is
returned. First fit, best fit, a bump pointer and block splitting all satisfy
it. The choice determines which addresses a structure occupies, never what a
\gls{cell} contains, so an allocation algorithm is not a semantic object.
\end{remark}

\term{pointer}{Pointer}
A pointer names a location in some region: which region, and how far into it.
Because region identifiers are ranks, a pointer says both where the bytes live
and which process is responsible for them --- an invariant the implementation
asserts on every operation.

\begin{definition}[Pointer]\label{def:pointer}
A \emph{pointer} is a global address $\gaddr{p}{a} \in \Addr$, well formed
when $a \in \Addrof{p}$. Write $\mathit{owner}(\pi) \defeq p$ and
$\mathit{off}(\pi) \defeq a$.
\end{definition}

% ===========================================================================
\subsection{Time}
\label{sec:gloss-time}

Four quantities are informally called ``time''. Separating them is the most
useful thing this document does.

\term{logical-time}{Logical time}
Logical time is simulation time: the timeline the simulated model advances
along, unrelated to wall-clock time or MPI progress. All processes share one
timeline and sit at their own points on it --- which is the entire purpose of
the design.

\begin{definition}[Logical time]\label{def:logical-time}
$\Time \defeq \Int$, ordered as usual. Times $\eff < 0$ are
\emph{pre-history}: no \gls{version} exists there, and a read directed at one
is answered from a caller-supplied default (\cref{def:get-or-default}).
\end{definition}

\term{origin-time}{Origin time}
The origin time of an operation is where the issuing process sits on the
timeline. It is supplied by the program, not inferred: supplying it is how a
process reports its position.

\term{delta}{Delta}
The delta is a signed displacement separating when an operation happens from
when its effect is visible. Writes carry $\delta \ge 0$, reads $\delta \le 0$.
It is the model's unit of latency, in simulated time.

\term{effective-time}{Effective time}
The effective time is the \gls{version} an operation actually touches: origin
plus delta. The memory model is indexed by it; origin time only tracks clocks.

\begin{definition}[Origin time, delta, effective time]\label{def:times}
An operation carries an \emph{origin time} $t \in \Time$ and a \emph{delta}
$\delta \in \Time$; its \emph{effective time} is $\eff \defeq t + \delta$.
Writes require $\delta \ge 0$, reads $\delta \le 0$.
\end{definition}

\term{local-clock}{Local clock}
A process's local clock is the furthest along the timeline it has declared
itself to be. It moves only forward, and only when the program issues an
operation with a larger origin time. It is what the process publishes to its
neighbours.

\begin{definition}[Local clock]\label{def:local-clock}
The \emph{local clock} $\clk{p} \in \Time$ of the process $\proc{p}$ is the
supremum of the origin times $\proc{p}$ has issued so far in the current
\gls{session}, or $0$ if none. It is nondecreasing.
\devmark{dev:clock-width}
\end{definition}

\term{known-time}{Known time}
A process cannot see its neighbours' clocks; it learns of them from messages.
The known time is its current belief about a neighbour --- always a
\emph{lower} bound, never an overestimate. That one-sidedness is what makes
the visibility rules safe.

\begin{definition}[Known time]\label{def:known-time}
For $p, q \in \Ranks$, the \emph{known time} $\kn{p}{q} \in \Time$ is
$\proc{p}$'s estimate of $\clk{q}$, initially $0$, updated only by
$\kn{p}{q} \mathrel{{:}{=}} \max(\kn{p}{q}, t)$ on receipt of a message from
$\proc{q}$ stamped $t$. It is nondecreasing.
\end{definition}

\begin{invariant}[Clock soundness]\label{inv:clock-soundness}
$\kn{p}{q} \le \clk{q}$ for all $p, q \in \Ranks$, always.
\end{invariant}

\term{drift}{Drift}
Drift is the clock skew between two processes. Tolerating it without rollback
is the point of the design; bounding it is what the \gls{comm-graph} does.

\begin{definition}[Drift]\label{def:drift}
$\drift{p}{q} \defeq \clk{p} - \clk{q}$.
\end{definition}

\term{mvt}{Minimum virtual time}
The oldest version a process still holds. A memory-management quantity, not a
semantic one: it says how far back the runtime is still paying to remember.

\begin{definition}[Minimum virtual time]\label{def:mvt}
$\mvt{p} \defeq \min \{\, \eff : \Vst{p}{\eff} \ne \Absent \,\}$, and
$\clk{p}$ when rank $p$ holds no version.
\end{definition}

% ===========================================================================
\subsection{Locality}
\label{sec:gloss-locality}

\term{comm-graph}{Communication graph}
The communication graph declares the shape of a simulation: which pairs of
ranks may ever exchange \glspl{operation}, and how long, in simulated time, an
interaction across such a pair takes to become visible. Supplied once per
\gls{session} and never changed, it is the source of every optimisation DSTM
performs --- the only reason a process can conclude that another can no longer
affect it.

\begin{definition}[Communication graph]\label{def:comm-graph}
A \emph{communication graph} is a triple $\cgraph = (\Ranks, \edges, w)$ in
which $\edges \subseteq \Ranks \times \Ranks$ is the \emph{channel} relation
and $w$ assigns each channel a weight, subject to
\begin{description}[leftmargin=2.9cm,style=nextline,font=\normalfont\bfseries]
  \item[\textnormal{(G1) Irreflexivity}] $(p,p) \notin \edges$.
  \item[\textnormal{(G2) Symmetry}] $(p,q) \in \edges \iff (q,p) \in \edges$,
    and $\lag{p}{q} = \lag{q}{p}$. \devmark{dev:asymmetric}
  \item[\textnormal{(G3) Positive integer weights}]
    $w : \edges \to \Nat_{>0}$. \devmark{dev:zero-distance}
\end{description}
\end{definition}

\term{lag}{Lag}
The lag of a channel is its weight: the number of simulated timesteps an
interaction across it takes to become visible. It is a \emph{time} separation
between two neighbours, nothing more.

\begin{convention}[Lag is not a distance]\label{conv:lag-not-metric}
$w$ is deliberately \emph{not} a metric and nothing below assumes it is. It is
defined only on channels, never between arbitrary ranks; no triangle
inequality is required, so a channel may cost more than a detour around it
without making the configuration ill formed. The only axioms are (G1)--(G3),
and only (G2) and (G3) are ever used.
\end{convention}

Of the three, (G3) carries the weight. Because $w$ is integer-valued it says
every channel costs at least one timestep, which makes every cycle in
$\cgraph$ carry positive weight --- and that is what makes the dependence
graph acyclic (\cref{thm:acyclic}) and the system deadlock-free
(\cref{thm:deadlock-freedom}). A configuration admitting a zero-weight channel
has none of DSTM's guarantees \devmark{dev:zero-distance}.

\term{reachability}{Adjacency and neighbourhood}
\alsoterm{neighbourhood}
Two ranks are \emph{neighbours} when a channel joins them directly.
Non-neighbours exchange no messages and neither constrains the other's
progress. This sparsity is what makes DSTM scale: synchronisation cost depends
on neighbourhood size, not on $\worldsize$.

\begin{definition}[Neighbourhood]\label{def:neighbourhood}
\[
  \nbr{p} \defeq \{\, q \in \Ranks : (p, q) \in \edges \,\},
  \qquad
  \diam_p \defeq \max_{q \in \nbr{p}} \lag{p}{q}
\]
with $\diam_p \defeq 0$ when $\nbr{p} = \emptyset$.
\end{definition}

\term{path-weight}{Path weight}
The path weight between two ranks is the least total lag of a chain of
channels joining them. Unlike the lag itself it is defined for every pair, and
it is the quantity that bounds how far apart two processes can drift
(\cref{prop:drift-bound}).

\begin{definition}[Path weight]\label{def:path-weight}
A \emph{path} from $p$ to $q$ is a sequence $p = z_0, \dots, z_k = q$ with
$(z_{m-1}, z_m) \in \edges$ for each $m$, of \emph{weight}
$\sum_{m} \lag{z_{m-1}}{z_m}$. The \emph{path weight}
$W : \Ranks \times \Ranks \to \Nat \cup \{\infty\}$ is the least weight of
such a path, with $\pw{p}{p} = 0$ and $\pw{p}{q} = \infty$ when no path
exists.
\end{definition}

\begin{remark}[Where a metric does belong]\label{rem:path-weight-metric}
$W$ is a metric --- a least-path weight always satisfies the triangle
inequality, by construction rather than by assumption --- whereas $w$ is not.
The distinction is exactly right: $w$ is data the application supplies and
must not be over-constrained, while $W$ is derived and is used only to bound
drift and memory. Note also that $\pw{p}{q} < \infty$ does \emph{not} make $p$
and $q$ neighbours: ranks at opposite corners of a mesh have a finite path
weight, share no channel, exchange no \glspl{heartbeat} and never appear in
each other's guards.
\end{remark}

\term{lookahead}{Lookahead}
Lookahead is the guarantee the \gls{comm-graph} buys: nothing $\proc{q}$ does
now can affect $\proc{p}$'s state earlier than $\lag{p}{q}$ into the future.
It is the same quantity conservative PDES calls lookahead \cite{chandy1979},
playing the same role --- what lets a process declare state final without
asking anyone.

\term{heartbeat}{Heartbeat}
A heartbeat carries only ``my clock has advanced to $t$''. Processes broadcast
one to their neighbours whenever their clock moves. It targets no address and
carries no data; it exists so neighbours can raise their known times and make
progress.

\begin{definition}[Heartbeat]\label{def:heartbeat}
A \emph{heartbeat} from $\proc{p}$ is a message $\langle p, t\rangle$ with
$t = \clk{p}$, delivered to every $q \in \nbr{p}$, whose sole effect is
$\kn{q}{p} \mathrel{{:}{=}} \max(\kn{q}{p}, t)$.
\end{definition}

Every message, not only a heartbeat, carries its sender's stamp and raises the
receiver's known time. A heartbeat is the degenerate message carrying nothing
else.

% ===========================================================================
\subsection{State}
\label{sec:gloss-state}

\term{cell}{Cell}
A cell is the smallest addressable unit of DSTM state, indexed by
\emph{three} coordinates: a rank, a time, and an offset. Time is part of the
location's identity, not metadata attached to it. Two writes to the same
offset at different effective times do not touch the same cell.

\begin{definition}[Cell]\label{def:cell}
$\Cells \defeq \Addr \times \Time$; we write $\cell{p}{\eff}{a}$ for the cell
at global address $\gaddr{p}{a}$ and \gls{effective-time} $\eff$
(\cref{def:cell-pgas}).
\end{definition}

\term{version}{Version}
\alsoterm{publication}
A version is the whole of one region at one \gls{effective-time} --- the cells
sharing a rank and a time. It is the unit the runtime allocates, tracks and
reclaims, created on demand by the first write targeting it, and filled with
zeroes: it does \emph{not} inherit its predecessor's contents.

\begin{definition}[Version]\label{def:version}
The \emph{version} $\ver{p}{\eff}$ is the family
$\{\, \cell{p}{\eff}{a} : a \in \Addrof{p} \,\}$ with content
$\Vbuf{p}{\eff} : \Addrof{p} \to \Bytes$ and status $\Vst{p}{\eff}$.
\end{definition}

\begin{convention}[Versions are independent, not incremental]
\label{conv:no-inherit}
$\Vbuf{p}{\eff}$ is initialised to all zeroes on creation, regardless of
$\Vbuf{p}{\eff'}$ for $\eff' < \eff$. Versions are independent snapshots
indexed by time, not deltas against an evolving state.
\devmark{dev:non-inheriting}
\end{convention}

This is the most consequential modelling decision in DSTM and the most easily
misread; \cref{sec:memory-model} takes it up.

\term{version-status}{Version status}
Every version passes through the same three stages. It is first
\emph{writable}, accepting writes and refusing reads; becomes \emph{readable}
once no further write can target it, freezing its contents; becomes
\emph{discardable} once no further read can target it, and is reclaimed. The
progression is one-way.

\begin{definition}[Status lattice]\label{def:status}
$\Status \defeq \{\Absent, \Wr, \Rd, \Dc\}$, totally ordered
$\Absent < \Wr < \Rd < \Dc$. A version's status is nondecreasing in it.
\end{definition}

\term{store}{Versioned store}
The total picture: which byte sits in every cell. This is what the memory
model is about and what an implementation must reproduce.

\begin{definition}[Versioned store]\label{def:store}
A \emph{versioned store} is a total function $\Store : \Cells \to \Bytes$.
\end{definition}

% ===========================================================================
\subsection{Operations}
\label{sec:gloss-ops}

\term{operation}{Operation}
\alsoterm{put}\alsoterm{get}
DSTM exposes exactly two data operations, a write and a read, both one-sided:
the issuing process drives it, and the target executes no application code on
its behalf. Each names a target rank, address, origin time and delta.

\begin{definition}[Operations]\label{def:operation}
An \emph{operation} is
\[
  \Put{p}{q}{t}{\delta}{a}{v}
  \qquad\text{or}\qquad
  \Get{p}{q}{t}{\delta}{a}{\ell},
\]
with $p \in \Ranks$ the \emph{issuer}, $q \in \Ranks$ the \emph{target},
$a \in \Addrof{q}$ the address, $v \in \Values$ the written value and
$\ell = |v| \ge 1$. We write $\opsrc{o}, \optgt{o}, \optime{o}, \opdelta{o}$
and $\opeff{o} = \optime{o} + \opdelta{o}$.
\end{definition}

\term{footprint}{Footprint}
The set of cells an operation touches. Because cells carry a time coordinate,
footprints do too --- which is why operations at different effective times can
never conflict however much their addresses overlap. \Cref{sec:sync} builds
the dependence analysis on this.

\begin{definition}[Footprint]\label{def:footprint}
For $o$ targeting rank $q$ at address $a$ with length $\ell$ and effective
time $\eff$,
\[
  \fp{o} \defeq \{\, \cell{q}{\eff}{a'} : a \le a' < a + \ell \,\}.
\]
$\fpw{o} \defeq \fp{o}$ if $o$ is a put and $\emptyset$ otherwise;
$\fpr{o}$ dually for gets. Operations are \emph{footprint-disjoint},
$o_1 \disjoint o_2$, when $\fp{o_1} \cap \fp{o_2} = \emptyset$.
\end{definition}

\term{issue}{Issue and completion}
\alsoterm{completion}
An operation is \emph{issued} when the program calls the API and
\emph{completes} when its effect is settled. The gap differs by operation: a
put returns immediately and completes later; a get does not return until it
completes. This is why a put carries no result and a get does.

\begin{definition}[Issue, completion]\label{def:issue-complete}
The \emph{issue} of $o$ is the instant its API call is made; its
\emph{completion} the instant its effect on the store (put) or its result
(get) is determined. Puts are \emph{asynchronous}, gets \emph{synchronous}.
\end{definition}

\term{history}{History and program order}
\alsoterm{program-order}
A history records everything a program did. Each process being
single-threaded, its own operations are totally ordered; operations of
different processes are not ordered a priori.

\begin{definition}[History, program order]\label{def:history}
A \emph{history} is a set $\Hist$ of operations with, for each $p \in \Ranks$,
a total order $\po$ on $\{\, o \in \Hist : \opsrc{o} = p \,\}$. Operations
issued by distinct ranks are unordered by $\po$.
\end{definition}

% ===========================================================================
\subsection{Coordination}
\label{sec:gloss-coord}

\term{quiescence}{Quiescence}
A system is quiescent when nothing is in flight and no process will issue
more. Detecting this is harder than detecting that every process has
\emph{finished issuing}, because an operation may still be travelling after
its issuer stopped; DSTM therefore detects the two conditions in sequence.

\begin{definition}[Quiescence]\label{def:quiescence}
A configuration is \emph{quiescent} when no message is in flight, no process
holds an incomplete operation, and none will issue another this
\gls{session}.
\end{definition}

\term{barrier}{Barrier}
DSTM's only collective inside the simulation loop, expected once at the end.
It establishes quiescence by circulating a token around the rank ring in three
rounds: the first establishes that every process has stopped issuing, the
second that every process has finished serving, the third releases everybody.
Two rounds are provably insufficient (\cref{thm:barrier-three-rounds}).

\term{termination}{Termination}
Quiescence plus teardown: after the barrier every pending-operation list is
empty and every region may be destroyed safely.

% ===========================================================================
\subsection{Derived structures}
\label{sec:gloss-derived}

The concepts above are the DSTM \emph{kernel}. Applications use a distributed
data structure layered on it; those layers add addressing discipline, not new
semantics.

\term{darray}{Distributed array (darray)}
\alsoterm{partition}
A one-dimensional array of fixed-size elements cut into contiguous slices, one
per rank. It supplies the map from a global element index to a
(rank, address) pair; reads and writes then go through the kernel unchanged.

\begin{definition}[Partition]\label{def:partition}
A \emph{partition} of $m$ elements is a vector
$(m_0, \dots, m_{\worldsize-1})$ of positive integers summing to $m$, with
\emph{offsets} $b_p \defeq \sum_{q < p} m_q$. The \emph{owner} of a global
index $g \in \{0, \dots, m-1\}$ is the unique rank $p$ with
$b_p \le g < b_p + m_p$, written $\mathit{own}(g)$.
\devmark{dev:darray-bounds}
\end{definition}

\begin{definition}[Distributed array]\label{def:darray}
A \emph{darray} is $(z, (m_p)_p, (\beta_p)_p)$ with element size $z \ge 1$, a
partition $(m_p)_p$, and $\beta_p \in \Addrof{p}$ the base of an
\gls{allocation} whose extent covers $[\beta_p, \beta_p + z\,m_p)$. Index $g$
resolves to
\[
  \mathit{res}(g) \defeq \bigl(\mathit{own}(g),\ \beta_{\mathit{own}(g)}
    + (g - b_{\mathit{own}(g)}) \cdot z\bigr).
\]
\end{definition}

The $\beta_p$ are chosen independently by each process and exchanged
collectively; the model must not assume they agree.

\term{dgraph}{Distributed graph (dgraph)}
A static directed graph with versioned state on vertices and edges. Topology
is frozen at construction; only state is versioned. An edge is owned by the
owner of its \emph{source} vertex, which puts the common ``consume my inbox
locally, produce into a neighbour's inbox'' pattern on the legal side of every
ownership boundary.

\term{topology-replica}{Topology replica}
Resolving an edge $(u,v)$ to an address requires the source rank's adjacency,
so each process keeps a copy of every neighbouring rank's topology. This
replica is static, unversioned and exchanged once at construction --- a second
class of memory, and the only one outside the versioned store, which is sound
because it is immutable.

% ===========================================================================
\subsection{Idioms}
\label{sec:gloss-idioms}

Two communication patterns recur and the API is shaped around them. They look
like opposites and are the same schedule (\cref{thm:pipeline-depth}).

\term{pull-model}{Pull model}
Each process publishes into its \emph{own} region at its current time
($\delta = 0$, local); consumers reach across into that region's \emph{past}
($\delta = -\lag{p}{q}$). Data sits still; readers reach back.

\term{push-model}{Push model}
Each process writes into its \emph{neighbours'} regions at a point in their
\emph{future} ($\delta = +\lag{p}{q}$); every process reads only its own
region at its current time ($\delta = 0$, local). Data moves; readers stay
home.

\term{mailbox}{Mailbox pattern}
Index the shared structure by \emph{(destination, source)}, so cell $(d,s)$ is
written only by $\proc{s}$ and read only by $\proc{d}$. Each cell then has
exactly one writer per effective time --- precisely the write-disjointness
determinism requires (\cref{res:write-disjoint}). Both benchmark applications
in \cite{dstm2026} use it.
